% !TEX root = ../main.tex

\chapter{绪论}

\section{研究背景与意义}

本科生毕业论文（设计）是本科人才培养中的重要环节，是考核学生掌握和运用所学基础理论、基本知识与基本技能解决实际问题能力的关键依据。随着计算机排版技术的发展，\LaTeX{} 因其卓越的公式编排能力、规范的交叉引用机制和跨平台一致性，已成为学术界与工程领域撰写高质量学位论文的首选工具。

西南财经大学（Southwestern University of Finance and Economics, SWUFE）作为教育部直属的“双一流”建设高校，在经济学、管理学、计算机科学与人工智能等多学科交叉领域取得了丰硕成果。为了帮助全校本科生摆脱 Word 排版中常见的编号错乱、格式不合规、公式模糊等痛点，特依照教务处发布的最新《西南财经大学本科生毕业论文（设计）撰写与印制规范》，研制本套专属于西财本科生的 \LaTeX{} 论文模板。

\section{西财本科论文撰写与印制规范概要}

根据教务处规范文件要求，本科生毕业论文的核心格式要点如下：

\subsection{基本版面与打印要求}
\begin{enumerate}
    \item \textbf{纸张与打印}：标准 A4 纸张，正反双面打印。
    \item \textbf{页边距}：全文页边距严格要求为\textbf{上 3 厘米、下 3 厘米、左 3 厘米、右 3 厘米}。
    \item \textbf{字数要求}：本科毕业论文正文字数一般在 \textbf{8000 字以上}，毕业设计正文字数一般在 \textbf{7000 字以上}。
    \item \textbf{页码编排}：
    \begin{itemize}
        \item 中文摘要页码：居中、小五号阿拉伯数字、起始页码为 1；
        \item 英文摘要页码：居中、小五号阿拉伯数字、起始页码为 1；
        \item 正文页码：居中、小五号阿拉伯数字、起始页码为 1。
    \end{itemize}
\end{enumerate}

\subsection{各级标题与正文字体字号}
\begin{table}[htbp]
\centering
\caption{西南财经大学本科毕业论文排版规范核心要求}
\label{tab:swufe-bachelor-rules}
\begin{tabularx}{\textwidth}{p{2.4cm}p{4.2cm}X}
\toprule
\textbf{内容部分} & \textbf{字体与字号} & \textbf{对齐与间距要求} \\
\midrule
一级标题 & 二号华文中宋加粗 & 居中，段前 3 行，段后 2 行，1.5 倍行距 \\
二级标题 & 小三号黑体 & 首行无缩进，段前 1 行，段后 1 行，1.5 倍行距 \\
三级标题 & 小四号黑体 & 首行无缩进，段前 1 行，段后 1 行，1.5 倍行距 \\
正文 & 小四号宋体 & 首行缩进 2 字符，两端对齐，段前 0，段后 0，1.5 倍行距 \\
中文摘要标题 & 二号华文中宋加粗 & 居中，段前 3 行，段后 2 行 \\
英文摘要标题 & 二号 Times New Roman 加粗 & 居中，段前 3 行，段后 2 行 \\
中文关键词 & 小四号宋体加粗 & “关键词：”加粗，分号分隔，末尾无标点 \\
英文关键词 & 小四号 Times New Roman 加粗 & “\textbf{Keywords:} ”加粗，分号分隔 \\
图/表标题 & 五号黑体加粗 & 表题在上方，图题在下方，段前段后空 0.5 行 \\
图/表内容 & 五号字体 & 单倍行距 \\
注释、脚注 & 小五号宋体或 Times New Roman & 每页独立编号（1, 2, 3……），正文上标 \\
参考文献 & 同一级标题 / 同正文 & 先中文后外文，起始序号为 [1] \\
致谢、附录 & 同一级标题 / 同正文 & 致谢内容限一页，置于附录之后 \\
\bottomrule
\end{tabularx}
\end{table}

\subsection{论文装订顺序}
论文胶装时必须严格按照以下顺序装订：
\begin{enumerate}
    \item 封面（本科毕业论文单封面，含校徽校名与手写下划线信息）
    \item 学术声明（本科生原创性及知识产权声明）
    \item 中文摘要及关键词
    \item 外文摘要及关键词
    \item 标题及目录（不需出现图、表目录）
    \item 正文（按章节组织）
    \item 参考文献
    \item 附录（如果有）
    \item 致谢（限一页）
    \item 封底（含“光华铁树”插图与竖排“经世济民 孜孜以求”）
\end{enumerate}

\section{模板快速配置指南}

本模板继承并升级了西南财经大学研究生模板的优质特性，支持\textbf{零依赖开箱即用}的 JSON 配置：

\begin{promptbox}
\textbf{推荐配置方式：} 直接编辑项目根目录下的 \texttt{config.json} 文件，填写个人姓名、学号、学院、专业、年级、导师等字段即可。无需修改宏包代码，无论在本地还是在 Overleaf 均可完美编译！
\end{promptbox}

如果不需要 \texttt{config.json}，直接将其删除，模板将自动回退至 \texttt{setup.tex} 中的 \texttt{\textbackslash sjtusetup} 原生参数。

\section{本文结构安排}

本文档既是本套本科毕业论文模板的技术示例文档，也是一份详细的排版指南。全文组织如下：
第一章介绍课题背景、规范要求及模板配置；
第二章详细示范数学公式、三线表、插图、代码块、算法与脚注的排版规范；
第三章展示毕业设计系统的架构设计与实现案例；
第四章对全文进行总结与展望。
